% Generated draft: requires mathematical and source review.
\documentclass[11pt]{article}
\usepackage[margin=25mm]{geometry}
\usepackage{fontspec}
\setmainfont{Latin Modern Roman}
\usepackage{amsmath,amssymb,mathtools,unicode-math}
\setmathfont{Latin Modern Math}
\usepackage{enumitem,needspace,xurl,hyperref}
\setlength{\parindent}{0pt}
\setlength{\parskip}{5pt}
\setlength{\emergencystretch}{3em}
\newcommand{\N}{\mathbb{N}}
\newcommand{\Z}{\mathbb{Z}}
\newcommand{\Q}{\mathbb{Q}}
\newcommand{\R}{\mathbb{R}}
\newcommand{\C}{\mathbb{C}}
\newcommand{\F}{\mathbb{F}}
\newcommand{\A}{\mathbb{A}}
\newcommand{\PP}{\mathbb{P}}
\newcommand{\E}{\mathbb{E}}
\newcommand{\eps}{\varepsilon}
\newcommand{\dd}{\,\mathrm{d}}
\begin{document}
\section*{UnsolvedMath ID: 30001713}
\textbf{Source problem code:} catalogue entry
\subsection{Definitions and mathematical statement}
The supplied record concerns a cluster algebra denoted by \(A(B,x)\), with an initial tuple written \(x=(x_1,\ldots,x_m)\). For nonnegative integer vectors \(a=(a_1,\ldots,a_m)\) and \(\mathbf d=(d_1,\ldots,d_m)\), write
\[
x^a=\prod_{i=1}^m x_i^{a_i},\qquad a\cdot\mathbf d=\sum_{i=1}^m a_i d_i.
\]
Thus \(a\cdot\mathbf d=0\) means that \(a_i d_i=0\) for every \(i\). Here nonnegative integer indexing is the interpretation used for the displayed exponents and dimension vectors.

The source uses \(J\), \(\operatorname{Rep}_{\mathbf d}(J)\), \(\operatorname{Irr}_{sr}\), and \(C_Z\) without supplying their definitions. In particular, it does not specify the algebra \(J\) or its relation to \(B\), the ground field for the representation spaces, the condition denoted by the subscript \(sr\), or the construction of the element \(C_Z\) associated with an indexed component \(Z\). These symbols retain their source-dependent meanings here.

\textbf{Conjecture as supplied.} For the intended cluster algebra \(A(B,x)\) and its intended associated data, the set
\[
G=\left\{x^a C_Z\;\middle|\;
 a,\mathbf d\in\mathbb N^m,\quad
 a\cdot\mathbf d=0,\quad
 Z\in\operatorname{Irr}_{sr}\bigl(\operatorname{Rep}_{\mathbf d}(J)\bigr)
\right\}
\]
is a basis of \(A(B,x)\).

The coefficient ring over which this is a basis is not specified. If that ring is denoted by \(R\), the basis assertion means that every displayed element belongs to \(A(B,x)\), and every \(f\in A(B,x)\) has a unique expression
\[
f=\sum_{g\in G}r_g g,\qquad r_g\in R,
\]
with only finitely many nonzero coefficients. This explains the intended algebraic basis property conditionally; it does not supply the missing choice of \(R\).

The permitted matrices \(B\), coefficient or frozen-variable conventions, and hypotheses on any potential are not supplied. Although the background calls the target a generic-potential formulation, it gives neither the potential nor the meaning of genericity. Consequently, the displayed conjecture is a faithful statement template whose full hypotheses and quantification require source clarification.
\subsection{Short English statement}
The conjecture asks whether the displayed monomials times the source-dependent elements \(C_Z\), indexed subject to disjoint exponent and dimension-vector supports, form a basis of the intended cluster algebra. The algebraic setup, indexing condition, construction of \(C_Z\), and coefficient ring remain unspecified. The supplied dated curation labels the problem partially solved and reports acyclic special cases; these claims have not been independently verified.
\subsection{Sources}
\begin{itemize}
\item Supplied local UnsolvedMath registry record, permanent ID \(30001713\), exact source problem code \texttt{catalogue entry}, title \emph{Generic Bases for Cluster Algebras}, in the Oberwolfach Reports Open Problems collection. No individual proposer is identified; the supplied proposed year is 2011.
\item Supplied report citation: \emph{Representation Theory of Quivers and Finite Dimensional Algebras} (2011), Oberwolfach Reports, pp.~523--608, DOI 10.4171/owr/2011/10. \url{https://doi.org/10.4171/owr/2011/10}. The record attributes the problem to workshop open-problem discussions documented by this report; no exact problem page or speaker is supplied.
\item Supplied literature source: Geiss--Leclerc--Schr\"oer, \url{https://arxiv.org/abs/1004.2781}. The record supplies no title, theorem number, or precise correspondence between this reference and the displayed formulation.
\item Supplied literature source: Oberwolfach Report 10/2011, \url{https://doi.org/10.4171/owr/2011/10}.
\item Literature assessment dated exactly \texttt{2026-08-21}: the supplied curation labels the problem partially solved and states that the generic-basis claim is proved for important acyclic cases, including the coefficient-free acyclic case, while the extracted generic-potential formulation is not established in full generality. This is supplied dated curation derived from the OWR material, not an independently verified literature assessment; its status and cited claims require independent verification.
\end{itemize}
\end{document}
